\scanpage{75}
puisse, pour tt tel $x$, choisir $y\in\overline{\{x\}}\cap Y$ tel que $k(y)$ soit \uline{fini}. \inserted{(c'est le cas si $\unsure{A/J|X/S}$ de type fini sur $\mathbf Z$).} Alors par l'argument précédent, on est ramené au cas d'un \deleted{anneau de val. \unsure{discrète|discret}} \inserted{localisé complété de} $X$ en $y$, \deleted{\unsure{donc|\illegible{1}} $\unsure{Y_0|U_0}$ \unsure{en $y$|$\subset y$}}, donc $X$ local complet, de pt fermé $y$, \unsure{avec|sur} $k(y)$ fini. Alors $X-Y$ est de Jacobson, donc tte composante connexe de $X-Y$ a un pt $z$ tel que $\dim\overline{\{z\}}=1$, \unsure{\uline{Soit}|\deleted{So}} $Z=\overline{\{z\}}$ et \unsure{son|$\overline Z$} l'adhérence dans $X$. On \unsure{se|en|nous} ramène, \deleted{par \unsure{un théorème|réduction}} quitte à remplacer $X$ par $Z$, au cas où $X$ local complet de dim 1, et en normalisant, au cas où $X$ est \inserted{le spectre d'} un anneau de valuation discrète \inserted{$V$} à corps résiduel fini, $Y$ pt fermé. L'hom envisagé correspond alors \unsure{à|:} un hom de modules sur $\pi=\operatorname{Gal}(\overline K/K)$, ext. de $\pi_0=\deleted{\unsure{\widehat{\mathbf Z}|\mathbf Z}}\begin{gathered}\operatorname{Gal}(\overline k/k)\\ \inserted{\Vert}\\ \inserted{\widehat{\mathbf Z}}\end{gathered}$ par grpe d'inertie $I$.

\deleted{On va \unsure{supposer} $T\simeq\mathbf G_{mk}^{\unsure{r|n}}$ (par extension résiduelle finie) et $T_\ell(\widetilde T)|U$ correspond \unsure{à|.}
\[
\deleted{T_\ell(\overline K^*)\simeq T_\ell(\overline k^*)}\quad (\widetilde{T_\ell(\mathbf G_{mk})}|U)^r\simeq T_\ell(\mathbf G_{mk})^r\simeq T_\ell(\overline k^*)^r,
\]
comme module sur $\deleted{\mathbf Z_\ell}\pi_0$. D'autre pt, \unsure{pour} $S$ \unsure{on|\illegible{1}} $T_\ell(S)$ \unsure{correspond|correspon}}

L'accouplement sur $U$, comme $X$ normal, provient d'un accouplement
\[
T_\ell(S)\times T_\ell(T)\mathbin{-}T_\ell(\mathbf G_{mk}).
\]
Donc dans ce cas on est ramené au

\uline{Lemme} Soit $k$ un corps fini, $\deleted{\unsure{G}}S,\unsure{\deleted{\unsure{G}}T'|\deleted{\unsure{G'}}T}$ deux grpes alg. commutatifs sur $k$, \deleted{\unsure{dont le second}} \inserted{ou $T$ de \unsure{t.|V.} mult.} \unsure{est un|\deleted{est un}} tore. Alors tt. accouplement
\[
T_\ell(S)\times T_\ell(T)\mathbin{-}T_\ell(\mathbf G_m)
\]
est identiquement nul.
\margin{74}
